\section{Conclusion\label{sec:conclusion}}

This study examines which implementation-level signals, including code text, static dependencies, folder organization, and member-level information, carry architectural information and how their representation and use affect code-to-architecture mapping on ten systems. We encode file information using lexical Word2Vec, structural Node2Vec, and contextual UniXcoder in zero-shot (\zs{}) and dependency-guided fine-tuned (\ft{}) variants, and evaluate these representations with MLP- and GNN-based predictors. Representations that encode dependencies, Node2Vec and \ft{}, outperform Word2Vec on most systems when evaluated with the MLP, whereas Word2Vec and \zs{}, which do not, perform comparably. Dependency-guided fine-tuning significantly improves mapping performance over \zs{} on most systems, although no representation is best on every system. Folder nodes produce the largest and most consistent gains within HGT variants, but largely recover information already available to a directory rule using the initial labels and do not consistently improve the files on which that rule fails. Member-node gains are less uniform and cannot generally be separated from changes in feature construction and placement.

The controlled experiments distinguish dependency information encoded in file representations from its additional use during prediction by keeping the input vectors fixed and switching direct dependency neighbors on and off for four predictors, under both training schedules and on all ten systems. Access to direct dependency neighbors improves performance across the evaluated representations for both GAT, which uses message passing, and the neighbor \mlp{}, which concatenates own-file features with mean-neighbor features, with larger gains for Word2Vec and \zs{} than for Node2Vec. These results show that dependency links provide useful information beyond lexical and pretrained code features, while gains with Node2Vec and \ft{} show that neighbor access can also complement representations that already encode dependencies. Whether neighborhood information helps depends on how it is processed: with Node2Vec vectors, GAT and \mlp{} improve macro-F1 on at least nine of ten systems under both schedules, whereas HGT lowers it on at least nine. Message passing over the directed, multi-relational dependency graph in HGT therefore benefits less than the simpler GAT and neighbor \mlp{}, and its median macro-F1 change is also negative with the other vectors. Removing direction and edge types from HGT removes most of this loss, but the two are changed together, so the experiments do not show which of the two accounts for it.

Dependencies carry architectural signal beyond file content, and a mapping model can use it in two ways: by encoding it in the file representation, as Node2Vec and dependency-guided fine-tuning do, or by giving the predictor access to direct neighbors. Both help, but the second depends on how the neighbors are processed. Evaluating a mapping model should therefore establish what its input already encodes, isolate neighbor access within a fixed predictor, and examine the files that depart from their directory assignment. These findings concern ten systems and the compared configurations, and do not rank the signals or the predictors in general.